% ======================================================================
%  Project Report: Tool Selection Robustness in LLM Agents
%  Compiled with pdflatex. Figures are optional: place PNG images in the
%  figures/ folder using the names referenced below, or the document
%  will still compile with placeholder boxes.
% ======================================================================
\documentclass[11pt,a4paper]{report}

\usepackage[margin=2.5cm]{geometry}
\usepackage{graphicx}
\usepackage{booktabs}
\usepackage{array}
\usepackage{enumitem}
\usepackage{caption}
\usepackage{float}
\usepackage[hidelinks]{hyperref}
\usepackage{xcolor}
\usepackage{titlesec}

\graphicspath{{figures/}}
\captionsetup{font=small,labelfont=bf}

\titleformat{\chapter}[display]
  {\normalfont\huge\bfseries}{\chaptertitlename\ \thechapter}{12pt}{\Huge}
\titlespacing*{\chapter}{0pt}{0pt}{18pt}

% Optional-figure helper: shows the image when present, otherwise a
% framed placeholder so the document always compiles.
\newcommand{\figimg}[3][0.85]{%
  \IfFileExists{#2}{%
    \begin{center}\includegraphics[width=#1\textwidth]{#2}\end{center}%
  }{%
    \begin{center}%
      \fbox{\parbox{#1\textwidth}{\centering\vspace{2.2cm}%
      [Place image here: #2]\vspace{2.2cm}}}%
    \end{center}%
  }%
  \caption{#3}%
}

\newcommand{\metric}[1]{\texttt{#1}}

\begin{document}

% ----------------------------------------------------------------------
\begin{titlepage}
  \centering
  {\LARGE\bfseries An Evaluation Framework for Tool Selection\\[2mm]
  Robustness in LLM Agents\par}
  \vspace{1.2cm}
  {\large Project Report\par}
  \vspace{0.6cm}
  {\large Prepared for research supervision and feedback\par}
  \vfill
  {\large Department of Computer Science and Engineering\par}
  \vspace{0.4cm}
  {\large \today\par}
\end{titlepage}

% ----------------------------------------------------------------------
\tableofcontents

% ======================================================================
\chapter{Introduction}
% ======================================================================

Large language models (LLMs) are increasingly deployed as \emph{agents}:
systems that plan a task and then call external tools, such as weather
services, calculators, or e-commerce APIs, to complete it. Before an agent
can execute anything, it must decide \emph{which tool} fits the user's
request. This decision, known as tool selection, is security critical. If
the agent picks the wrong tool, or a tool that was crafted with harmful
intent, the consequences can include data leaks, unauthorized actions, or
simply incorrect answers.

Two properties of the real world make this problem hard. First, the space
of available tools is large and constantly growing, so agents narrow the
search using a retrieval step before making a final choice with an LLM.
Second, tool repositories are open. Anyone can publish a tool together
with a short natural language description of what it does, and those
descriptions are later fed directly into the LLM's context. A description
written to exploit this channel can carry hidden instructions that
manipulate the selection process itself.

This report describes an evaluation framework we built to study these
questions in a controlled setting. The framework implements the complete
tool selection pipeline, provides a way to inject crafted tool documents
into the library, measures how strongly those documents influence
selection, and evaluates a set of detection methods that aim to flag
suspicious documents before they are ever selected. Everything runs
inside a synthetic benchmark environment, and the framework ships with an
interactive dashboard so that every measurement can be inspected and
reproduced visually.

The rest of the report is organized as follows. Chapter 2 gives the
background on LLM agents and tool selection. Chapter 3 defines the threat
model we assume. Chapter 4 describes the system architecture. Chapter 5
explains the evaluation methodology and metrics. Chapters 6 and 7 cover
the two studies the framework supports: the injection study and the
detection study. Chapter 8 summarizes the implementation and the
dashboard. Chapter 9 presents illustrative results obtained with the
current synthetic setup. Chapters 10 and 11 discuss limitations and
conclusions.

% ======================================================================
\chapter{Background}
% ======================================================================

\section{LLM Agents and Tools}

An LLM agent is an application that wraps a language model with the
ability to act on the world. Typical capabilities include reading a
knowledge base, browsing a website, and, most importantly for this work,
calling external tools. Each tool is an API with a defined function. What
the model sees of the tool is not the implementation but a short
\emph{tool document}: a name plus a natural language description of what
the tool does.

\section{The Two Step Selection Pipeline}

Modern agent frameworks rarely ask the LLM to pick directly from a
catalog of thousands of tools. Instead they decompose selection into two
steps, which we adopt in our framework as well.

\begin{enumerate}
  \item \textbf{Retrieval.} A lightweight embedding model maps the user's
  query and every tool document into vector space. The documents most
  similar to the query are kept, producing a short candidate list of size
  $k$ (the top-$k$ set).
  \item \textbf{Selection.} The user's query and the $k$ candidate
  documents are placed into a structured prompt for the LLM, together
  with strict instructions to choose exactly one tool and to answer in a
  fixed JSON format.
\end{enumerate}

Figure~\ref{fig:pipeline} illustrates the pipeline. The two steps have
different failure modes, and any security analysis must consider both: a
document must first \emph{get retrieved} and then \emph{get selected}.

\begin{figure}[H]
  \centering
  \figimg{pipeline.png}{The two step tool selection pipeline: the query
  is embedded, the top-$k$ tool documents are retrieved by similarity,
  and the LLM selects exactly one tool.}
  \label{fig:pipeline}
\end{figure}

\section{Why Tool Documents Are a Risk}

The descriptions inside tool documents are untrusted data by nature. They
are written by third parties, fetched from open repositories, and then
concatenated into the same prompt that carries the system's own
instructions. Because LLMs cannot perfectly separate instructions from
data, a description can be written so that the model treats part of it as
an order: for example, a sentence instructing the model to always prefer
that particular tool for a certain kind of request. This phenomenon is
generally known as prompt injection, and the tool document channel is one
of its most direct instantiations.

% ======================================================================
\chapter{Threat Model}
% ======================================================================

We model a platform where third parties can publish tools, mirroring
public tool hubs used by real agent frameworks. The evaluator plays the
role of such a third party and studies how much influence one published
document can exert.

\section{Objective}

The objective is to answer a measurable question: given a target task,
such as answering weather queries, can a single crafted tool document,
inserted into an otherwise benign library, redirect the agent's selection
towards the crafted tool across a wide range of natural user phrasings of
that task?

\section{Assumptions}

We adopt a conservative, black box setting, which we call the no box
setting, for the document author:

\begin{itemize}
  \item The author knows the target task, but does not know the exact
  queries real users will type.
  \item The author cannot read the victim's tool library, cannot query
  the victim's retriever or LLM, and does not know the top-$k$ setting.
  \item The author can only publish one document: a name and a
  description.
